\documentclass[11pt]{article}

\usepackage[margin=1in]{geometry}
\usepackage{amsmath,amssymb,amsthm}
\usepackage{mathtools}
\usepackage{enumitem}
\usepackage{xcolor}
\usepackage{hyperref}
\hypersetup{colorlinks=true,linkcolor=blue!55!black,citecolor=blue!55!black,urlcolor=blue!55!black}

\newtheorem{theorem}{Theorem}
\newtheorem{lemma}[theorem]{Lemma}
\newtheorem{proposition}[theorem]{Proposition}
\newtheorem{corollary}[theorem]{Corollary}
\theoremstyle{definition}
\newtheorem{definition}[theorem]{Definition}
\newtheorem{remark}[theorem]{Remark}

\DeclareMathOperator{\PP}{\mathbb{P}}
\DeclareMathOperator{\EE}{\mathbb{E}}
\newcommand{\N}{\mathbb{N}}
\newcommand{\eopt}{E_{\mathrm{opt}}}
\newcommand{\Fg}{F_G}
\newcommand{\bigoh}{O}
\newcommand{\half}{\tfrac12}

\title{The exact expected score of card guessing\\ after an $m$-shelf shuffle}
\author{Matt Watts\thanks{Independent researcher, \texttt{matt.ryan.watts@gmail.com}.
This work was carried out with substantial AI assistance (see the Acknowledgments);
all results are machine-verified by exact-arithmetic enumeration.}}
\date{July 2026\\[2pt]\emph{Working draft}}

\begin{document}
\maketitle

\begin{abstract}
A deck of $n$ distinct cards is shuffled by a single pass of a Diaconis--Fulman--Holmes
$m$-shelf shuffling machine and then dealt face up one card at a time; before each card a
player guesses its identity and is then shown the card (\emph{complete feedback}). We
determine the exact expected number of correct guesses of the optimal player, up to an
exponentially small remainder, for every number of shelves $m$. Writing $H_k=\sum_{j\le
k}1/j$ and $H_k^{(2)}=\sum_{j\le k}1/j^{2}$, the score of the Diaconis--Fulman--Holmes
strategy $G$ is
\[
  \Fg(n,m)\;=\;\frac{H_{2m}}{2m}\,n\;+\;\Bigl(\tfrac32-\tfrac{1}{4m}-H_{2m}^{(2)}\Bigr)
  \;+\;\bigoh\!\left(n^2\bigl(1-\tfrac1m\bigr)^{\!n}\right).
\]
The leading coefficient $H_{2m}/2m$ is the term conjectured by Clay (2025), whose general-$m$
proof he left open; we prove it, together with the exact constant term $\tfrac32-\tfrac1{4m}
-H_{2m}^{(2)}$ (limit $\tfrac32-\pi^2/6$) and the exponential decay rate $1-1/m$, by a direct
analysis of the shuffle's block structure that never forms the ``$m$-shelf transition
matrix'' Clay identified as the obstruction. Strategy $G$ is verified to be optimal exactly
for all $n\le 9,\ m\le 10$ and numerically to within $0.01$ for $m\le 40$ at deck scale; under
that (Diaconis--Fulman--Holmes / Clay) optimality the displayed expansion is the optimal
expected score.
\end{abstract}

\section{Introduction}

\subsection*{The problem}
Consider a deck of $n$ distinct cards, labelled $1,\dots,n$ by their position in the deck
before shuffling. The deck is shuffled by one pass of an $m$-shelf shuffling machine (defined
in \S\ref{sec:model}) and then dealt face up, one card at a time. Before each card is turned
over, a player names a card; the card is then revealed, and the player scores a point if the
guess was correct. This is card guessing \emph{with complete feedback}: the player sees every
card immediately after guessing it. The player wants to maximise the expected number of
correct guesses.

Because the player's guesses do not influence the order in which the cards appear, the optimal
strategy is the greedy Bayesian one: at each step, guess the card of maximal posterior
probability given the cards seen so far. Let $\eopt(n,m)$ denote the resulting optimal expected
score. Under a \emph{uniformly} random shuffle every strategy scores $H_n=\sum_{j=1}^n 1/j
\approx 4.5$ at $n=52$; an $m$-shelf shuffle is far from uniform, so a player who understands
the machine does strictly better.

\subsection*{Background}
Shelf shuffling machines and the associated guessing problem were introduced and analysed by
Diaconis, Fulman and Holmes~\cite{DFH}, who studied a commercial machine with $m=10$ shelves.
They proposed a simple, \emph{shelf-count-independent} guessing strategy $G$ --- guess the
successor of the last card shown, reversing direction near the top of the range --- and
conjectured it optimal, reporting an expected score of about $9.3$ out of $52$ at $m=10$.

The single-shelf case $m=1$ was settled by Clay~\cite{Clay}: the optimal strategy is exactly
$G$ (\cite[Thm.~1.4]{Clay}) and the optimal expected score is $\tfrac34 n$
(\cite[Thm.~1.5]{Clay}). For general $m$, Clay's Conjecture~3 asserts that $G$ is optimal (in a
high-probability sense, for $n/m$ not too small) with expected score
\begin{equation}\label{eq:clay}
  F_G(n,m)\ \approx\ \frac{n}{2m}\,H_{2m},
\end{equation}
reducing to $\tfrac34 n$ at $m=1$. Clay states that ``the transition matrix for an arbitrary
number of shelves is an open problem'' and closes with the hope ``to make this argument
precise'' in future work. A cluster of very recent papers \cite{Tripathi,Kuba,CKT} refines the
\emph{single-shelf} theory --- including the eigenvalues of the single-shelf \emph{position}
matrix, which governs the different, no-feedback game --- so the multi-shelf ($m\ge2$)
complete-feedback problem, on which \eqref{eq:clay} sits, has remained open.

\subsection*{Results}
We prove \eqref{eq:clay} for all $m$, as the leading term of a sharper three-term expansion.

\begin{theorem}\label{thm:main}
Let $H_k=\sum_{j=1}^k \frac1j$ and $H_k^{(2)}=\sum_{j=1}^k \frac1{j^2}$. For the $m$-shelf
shuffle with complete feedback, the expected score of the Diaconis--Fulman--Holmes strategy
$G$ is
\begin{equation}\label{eq:main}
  \Fg(n,m)\;=\;\frac{H_{2m}}{2m}\,n\;+\;b(m)\;+\;\bigoh\!\left(n^2\bigl(1-\tfrac1m\bigr)^{n}\right),
  \qquad
  b(m)=\frac32-\frac1{4m}-H_{2m}^{(2)}.
\end{equation}
The intercept satisfies $b(1)=0$ and $b(m)\to \tfrac32-\pi^2/6\approx-0.14493$ as $m\to\infty$;
at $m=1$ the remainder vanishes and \eqref{eq:main} reads $\Fg(n,1)=\tfrac34 n$, recovering
Clay's Theorem~1.5.
\end{theorem}

Three facts are proved here, each new for $m\ge 2$: the \emph{slope} $H_{2m}/2m$
(\S\ref{sec:slope}), the \emph{fade rate} $1-1/m$ (\S\ref{sec:fade}), and the exact
\emph{intercept} $b(m)$ (\S\ref{sec:intercept}), the last being the main contribution. The
method is elementary and probabilistic: we recast the shuffle as an assignment of independent
uniform labels (\S\ref{sec:model}), extract a single exchangeability lemma
(\S\ref{sec:lemma}), and sum a per-position hit probability over the deck. It never forms the
transition operator Clay named as the obstruction.

The relationship of \eqref{eq:main} to the \emph{optimal} score $\eopt$ is the content of
Clay's other claim --- that $G$ is optimal. We do not settle that in closed form for all $m$;
we verify it (\S\ref{sec:opt}): $\eopt=\Fg$ holds \emph{exactly} for all $n\le 9,\ m\le 10$,
and to within $0.01$ for all $m\le 40$ at $n\in\{52,104\}$, and we show that $G$ is the
optimal guess conditional on the true block-parse (Lemma~\ref{lem:key}). Under the
Diaconis--Fulman--Holmes/Clay optimality conjecture, \eqref{eq:main} is therefore the optimal
expected score. Every displayed identity in this paper has been checked by exact rational
enumeration; see \S\ref{sec:verify}.

\section{The \texorpdfstring{$m$}{m}-shelf shuffle as a block model}\label{sec:model}

We use the following clean description of the $m$-shelf shuffle, equivalent to the machine of
\cite{DFH} (each of the $m$ shelves receives cards on its top or its bottom; a shelf together
with a top/bottom orientation is one of $2m$ ``lanes'').

\begin{definition}\label{def:shuffle}
Assign to each card $c\in\{1,\dots,n\}$ an independent \emph{label}
$L_c$ uniform on $\{0,1,\dots,2m-1\}$. Order the cards by the key
\[
  \kappa_c=\begin{cases}(L_c,\,+c) & L_c\ \text{even},\\[2pt] (L_c,\,-c) & L_c\ \text{odd},\end{cases}
\]
increasingly in lexicographic order. The resulting sequence $o_1,\dots,o_n$ is the shuffled
deck.
\end{definition}

Because the key is lexicographic with the label first, the output is the concatenation of the
$2m$ \emph{blocks}
\[
  B_\ell=\{c: L_c=\ell\},\qquad \ell=0,1,\dots,2m-1,
\]
each written \emph{internally} in increasing order of value when $\ell$ is even and decreasing
when $\ell$ is odd:
\[
  \text{output}\;=\;\underbrace{B_0}_{\uparrow}\ \underbrace{B_1}_{\downarrow}\ \underbrace{B_2}_{\uparrow}\ \cdots\ \underbrace{B_{2m-1}}_{\downarrow}.
\]
Thus the output is $2m$ consecutive monotone runs, ascending and descending alternately, and
has at most $m-1$ interior local minima --- the class law of \cite[Thm.~3.1]{DFH}. At $m=1$
there are two blocks, $B_0\!\uparrow$ then $B_1\!\downarrow$.

We call the assignment $(L_c)_c$ the \emph{parse}: it records which block each card lies in.
The player observes only $o_1,\dots,o_n$, never the labels.

\section{The guessing game, the strategy \texorpdfstring{$G$}{G}, and the key lemma}\label{sec:lemma}

At step $t$ the player has seen $o_1,\dots,o_{t-1}$ and guesses a value; the optimal
(posterior-maximising) guess scores with probability $h_t:=\max_c \PP(o_t=c\mid
o_1,\dots,o_{t-1})$, and $\eopt(n,m)=\EE\sum_{t=1}^n h_t$.

\begin{definition}[strategy $G$]\label{def:G}
$G$ guesses $o_1$ to be the value $1$. For $t\ge2$, let $v=o_{t-1}$. If the last observed step
was ascending ($o_{t-1}>o_{t-2}$, or $t=2$) guess the smallest unrevealed value exceeding $v$;
if descending, the largest unrevealed value below $v$; and if the indicated side is empty,
take the nearest unrevealed value on the other side. $G$ uses only the observed order.
\end{definition}

This is the rule of \cite{DFH}: continue in the current direction, reversing at the ends. Let
$\Fg(n,m)$ be its expected score.

\begin{lemma}[parse-conditional exchangeability]\label{lem:key}
Fix a revealed prefix $o_1,\dots,o_s$ and condition on its true parse. Suppose its maximal
monotone runs are the nonempty parts of $B_0,\dots,B_\ell$; that blocks $0,\dots,\ell-1$ are
fully revealed; and that the current run (say ascending, $\ell$ even) ends at $o_s=v$. Then the
labels of the unrevealed cards are independent, with
\[
  L_c\sim \mathrm{Unif}\{\ell,\ell+1,\dots,2m-1\}\ \ (c>v),\qquad
  L_c\sim \mathrm{Unif}\{\ell+1,\dots,2m-1\}\ \ (c<v).
\]
Consequently the next card is the minimum-key unrevealed card; writing $w_1<w_2<\cdots$ for the
unrevealed values above $v$, one has $\PP(o_{s+1}=w_j)=\frac1r\bigl(1-\frac1r\bigr)^{j-1}$ with
$r=2m-\ell$, so the posterior-maximising guess is $w_1$ --- exactly $G$'s guess --- and its hit
probability is largest at $j=1$. (The descending case is the mirror image.)
\end{lemma}

\begin{proof}
The prior $\PP$ is a product measure over the cards' labels. The event ``the length-$s$ output
prefix equals $o_1\cdots o_s$'' is the intersection of per-card events: each revealed card's
label equals its parsed block index; and each unrevealed card $c$ has key exceeding
$\kappa_v=(\ell,v)$, i.e.\ $L_c>\ell$, or $L_c=\ell$ with $c>v$. A product prior intersected
with per-card constraints has a product posterior, uniform on each card's surviving label set.
For $c>v$ the label $\ell$ survives (the tie $L_c=\ell$ is broken by $c>v$), giving
$\{\ell,\dots,2m-1\}$; for $c<v$ it does not (a label-$\ell$ card below $v$ would have been
revealed before $v$), giving $\{\ell+1,\dots,2m-1\}$. The next output card is the one of
minimum key; among the unrevealed cards those below $v$ carry labels $\ge\ell+1$, hence larger
keys than any surviving label-$\ell$ card above $v$, so $o_{s+1}=w_j$ iff $L_{w_j}=\ell$ while
$L_{w_1},\dots,L_{w_{j-1}}>\ell$, which is geometric with success probability
$1/r$.\end{proof}

Two features of Lemma~\ref{lem:key} drive everything below. First, the \emph{bulk hit} is
$1/r=1/(2m-\ell)$: conditional on being in block $\ell$ with unrevealed cards on the
continuation side, $G$ guesses correctly with probability $1/(2m-\ell)$ up to a finite-size
correction. Second, $G$'s guess is the posterior-maximiser \emph{given the parse}; the player
does not know the parse, and it is precisely the resulting parse-uncertainty --- concentrated
at block boundaries and at the ends of the value range --- that produces the constant term
$b(m)$.

\section{The leading term}\label{sec:slope}

\begin{proposition}\label{prop:slope}
$\Fg(n,m)=\dfrac{H_{2m}}{2m}\,n+\bigoh(1)$, and the same holds for $\eopt$.
\end{proposition}

\begin{proof}
Call a guess position \emph{interior} if its prefix has an unambiguous parse and an unrevealed
value on the continuation side. By Lemma~\ref{lem:key} an interior position in block $\ell$ has
hit $1/(2m-\ell)+\text{(finite-size correction)}$, and the corrections sum to $\bigoh(1)$ (they
are handled exactly in \S\ref{sec:intercept}). The non-interior positions are $\bigoh(1)$ in
expectation: block transitions number at most $2m-1$; the value-range endgame contributes
$\bigoh(m)$; and a parse-ambiguous prefix requires an interior block to be momentarily empty,
an event of probability $\le 2m(1-\tfrac1{2m})^n$, so the expected number of such positions is
$\bigoh(n(1-\tfrac1{2m})^n)=\bigoh(1)$.

Each block has expected size $\EE|B_\ell|=n/2m$ (uniform labels), and contributes
$|B_\ell|-\bigoh(1)$ interior positions with block index $\ell$. Summing the bulk hits,
\[
  \Fg(n,m)=\sum_{\ell=0}^{2m-1}\frac{1}{2m-\ell}\Bigl(\frac{n}{2m}-\bigoh(1)\Bigr)+\bigoh(1)
  =\frac{n}{2m}\sum_{\ell=0}^{2m-1}\frac1{2m-\ell}+\bigoh(1)
  =\frac{H_{2m}}{2m}\,n+\bigoh(1).
\]
The lower bound is $G$'s explicit score; the upper bound for $\eopt$ follows because the
observer's actual hit is a mixture over parses of parse-conditional hits, each $\le
1/(2m-\ell)$ up to the same $\bigoh(1)$ endgame, and mixing cannot raise the sum above
$(H_{2m}/2m)n+\bigoh(1)$ by the tower property.
\end{proof}

The identity $H_{2m}=\sum_{\ell=0}^{2m-1}\tfrac{1}{2m-\ell}$ demystifies Clay's harmonic
number: it is the sum over blocks of $1/(\#\text{still-live blocks})$, the observer's advantage
from having watched exhausted blocks disappear. At $m=1$ the two blocks give hits $\tfrac12$
and $1$, slope $\tfrac34$: Clay's Theorem~1.5.

\section{The fade rate}\label{sec:fade}

\begin{proposition}\label{prop:fade}
The remainder in Theorem~\ref{thm:main} is $\Theta\!\left(n^2(1-\tfrac1m)^n\right)$: the
dominant subdominant eigenvalue of the guessing dynamics is exactly $1-1/m$, with algebraic
multiplicity $3$, so its Jordan structure carries a degree-$2$ polynomial prefactor.
\end{proposition}

\begin{proof}[Proof sketch]
Consider a contiguous ascending prefix $(1,2,\dots,k)$ (block $0$). The output begins with it
iff every unrevealed card $c>k$ has key exceeding $\kappa_k$; conditioning on the true block
$\ell=L_k$ of the last card, each such card survives \emph{independently} with probability
$\rho_\ell=(2m-\ell)/2m$ ($\ell$ even) or $(2m-1-\ell)/2m$ ($\ell$ odd). Hence
$\PP(\text{prefix},\,L_k=\ell)=K_\ell\,\rho_\ell^{\,n-k}$ with $K_\ell$ independent of $n$. The
observer's excess score over the bulk value $1/2m$ is carried by the competing hypotheses
$\ell\neq0$, whose posterior weight decays like $(\rho_\ell/\rho_0)^n$. The dominant competitor
is $\ell\in\{1,2\}$, both with $\rho_\ell/\rho_0=(2m-2)/2m=1-1/m$ --- a same-direction block
skip $\ell\to\ell+2$, the source of the factor $2$ in the exponent, giving the exponential
rate $1-1/m$. This rate is attained with algebraic multiplicity $3$ --- an exact property of
the guessing operator's spectrum, read off the exact-rational recurrence in
\S\ref{sec:verify} --- whose Jordan structure multiplies the geometric by a degree-$2$
polynomial, so the correction is $\Theta(n^2(1-\tfrac1m)^n)$. At $m=1$ there is no competing
hypothesis and the remainder is exactly $0$: the expansion is exact for $n\ge2$.
\end{proof}

\section{The exact intercept}\label{sec:intercept}

This is the main new computation. We evaluate the constant term of $\Fg$ in closed form.

\subsection{Reduction to a finite-size excess}
Index guess positions $t=2,\dots,n$ by the last-seen card $v=o_{t-1}$, in block $\ell_t=L_v$.
Assign to position $t$ the \emph{reference} value $1/(2m-\ell_t)$ (its bulk hit). Summing the
reference over all $n$ output positions gives $\sum_{c}1/(2m-L_c)$, with expectation
$n\,H_{2m}/2m=(\text{slope})\cdot n$ exactly; the last output card carries no follow-up guess,
and its label is $2m-1$ with probability $1-\bigoh((1-\tfrac1{2m})^n)$, contributing a
reference $1/(2m-(2m-1))=1$. The first guess ($v$ undefined) hits with probability
$\PP(o_1=1)=\PP(L_1=0)=1/2m$. Writing
\[
  S(m):=\lim_{n\to\infty}\sum_{t=2}^{n}\EE\!\left[h_t^{G}-\frac{1}{2m-\ell_t}\right],
\]
where $h^G_t$ is $G$'s hit at position $t$, we obtain
\begin{equation}\label{eq:bS}
  b(m)=-1+\frac1{2m}+S(m).
\end{equation}
It remains to evaluate $S(m)$, the total finite-size excess of $G$'s hit over its bulk value.

\subsection{Four bins}
$G$'s direction at a position is the last \emph{observed} step's direction; the block parity
of $v$ is hidden. Split each position by two binary attributes:
\begin{itemize}[topsep=2pt,itemsep=1pt]
  \item \textbf{on-foot / off-foot:} does $G$'s direction match the parity of $v$'s block
        ($\uparrow$ for even, $\downarrow$ for odd)? A position is off-foot exactly when $v$ is
        the first card of a block entered against that block's monotone direction; there $G$
        guesses on the wrong side and (a.s.\ in the limit) misses.
  \item \textbf{C / E:} is $G$'s chosen side non-empty (a genuine continuation, $C$) or empty
        (a peak/valley flip, $E$)?
\end{itemize}
Only the on-foot continuation bin $\mathsf{onC}$ contains $\Theta(n)$ positions; the other
three are $\bigoh(1)$, bounded by $\sim2m$ distinguished cards per shuffle (block-firsts,
block-maxima, block-minima). We show
\begin{equation}\label{eq:pieces}
  \mathsf{onC}=\frac{H_{2m}-1}{2},\quad
  \mathsf{offC}=1-\frac1{4m}-\frac{H_{2m}}{2},\quad
  \mathsf{onE}=2-\frac1{2m}-H_{2m}^{(2)},\quad
  \mathsf{offE}\to0,
\end{equation}
whose sum is $S(m)=\tfrac52-\tfrac{3}{4m}-H_{2m}^{(2)}$; with \eqref{eq:bS} this gives
$b(m)=\tfrac32-\tfrac1{4m}-H_{2m}^{(2)}$, proving Theorem~\ref{thm:main}.

\subsection{The interior bin \texorpdfstring{$\mathsf{onC}$}{onC}}
On an on-foot continuation in block $\ell$ (say ascending) with $A$ unrevealed cards above $v$
(including the guess target) and $B$ below, Lemma~\ref{lem:key} and a count of the independent
uniform labels give the exact hit
\begin{equation}\label{eq:hit}
  \PP(o_{t}=w_1)=\frac1r\sum_{d=0}^{r-1}
  \Bigl(1-\tfrac{d+(d\bmod 2)}{r}\Bigr)^{\!A-1}\Bigl(1-\tfrac{d-(d\bmod 2)}{r-1}\Bigr)^{\!B},
  \qquad r=2m-\ell,
\end{equation}
the $d=0$ term being $1$ (the bulk hit $1/r$) and $d\ge1$ the finite-size excess. Now sum the
excess over all positions. Given $L_v=\ell$, the unrevealed counts are independent binomials
$A\sim\mathrm{Bin}(n-v,\,r/2m)$ and $B\sim\mathrm{Bin}(v-1,\,(r-1)/2m)$, and $\PP(L_v=\ell)=1/2m$.
For the $d$-th term put $a_d=d+(d\bmod2)$ (even, $\ge2$) and $c_d=d-(d\bmod2)$; a direct
computation of the binomial generating functions gives, summed over $v=1,\dots,n$,
\[
  \frac{1}{2m}\cdot\frac1r\cdot\frac{1}{1-a_d/r}\sum_{v=1}^{n}
  \Bigl(1-\tfrac{c_d}{2m}\Bigr)^{v-1}\!\Bigl[\Bigl(1-\tfrac{a_d}{2m}\Bigr)^{n-v}-\Bigl(\tfrac{\ell}{2m}\Bigr)^{n-v}\Bigr].
\]
As $n\to\infty$ this vanishes whenever $c_d>0$ (the factor $(1-c_d/2m)^{v-1}<1$ kills the
geometric sum), i.e.\ for all $d\ge2$. Only $d=1$ survives ($c_1=0$, so the $B$-factor is
$1$), and there $a_1=2$, giving the limit
\[
  \frac{1}{2m\,r}\cdot\frac{1}{1-2/r}\Bigl(m-\frac{2m}{r}\Bigr)=\frac{1}{2m\,r}\cdot m=\frac{1}{2r}.
\]
Summing over the blocks $\ell=0,\dots,2m-2$ (i.e.\ $r=2,\dots,2m$; the top block $r=1$ has no
$d\ge1$ term),
\[
  \mathsf{onC}=\sum_{r=2}^{2m}\frac{1}{2r}=\frac{H_{2m}-1}{2}.
\]
The descending blocks give the same sum by the value-reflection symmetry $v\mapsto n{+}1{-}v$,
$\ell\mapsto 2m{-}1{-}\ell$; they are already included in the range of $\ell$.

\subsection{The off-foot bin \texorpdfstring{$\mathsf{offC}$}{offC}}
Off-foot positions are the first cards of blocks entered against their monotone direction, of
which there are at most $2m-1$. Consider block $\ell\in\{1,\dots,2m-1\}$ and its first output
card $F_\ell$. Take $\ell$ odd (descending); $F_\ell$ is the maximum of $B_\ell$, and it is
off-foot exactly when it is reached by an upward step, i.e.\ $F_\ell$ exceeds the last card of
the preceding block, i.e.\ $\max B_\ell>\max B_{\ell-1}$. Then $G$ guesses upward, but the next
card is the second card of $B_\ell$, which lies below $F_\ell$; so $G$ misses (hit $0$),
provided $B_\ell$ has a second card and there is an unrevealed value above $F_\ell$ (else the
position is in $\mathsf{offE}$). Its contribution to $S(m)$ is therefore
$-\,\PP(\text{off-foot, live continuation})\cdot\frac{1}{2m-\ell}$.

The probability is a first-occurrence count. Scanning the values from the top downward, the
event is: the first label in $\{\ell-1,\ell\}$ to appear is $\ell$ (probability $\tfrac12$),
and at least one label in $\{\ell+1,\dots,2m-1\}$ appears before it (else there is no
unrevealed value above $F_\ell$). Conditioning on all $2m$ blocks being nonempty (which holds
with probability $1-\bigoh((1-\tfrac1{2m})^n)$, and is exact in the limit), the first-occurrence
order of the labels $\{\ell-1,\ell\}\cup\{\ell+1,\dots,2m-1\}$ is uniform, so
\[
  \PP(\text{off-foot, live})=\tfrac12-\frac{1}{2m+1-\ell}
\]
(the subtracted term is the probability that $\ell$ precedes all of $\ell-1,\ell+1,\dots,2m-1$).
The even (ascending) blocks give the identical expression by reflection. Summing over
$\ell=1,\dots,2m-1$ with $k=2m-\ell$,
\[
  \mathsf{offC}=-\sum_{k=1}^{2m-1}\frac1k\Bigl(\frac12-\frac1{k+1}\Bigr)
  =-\frac12 H_{2m-1}+\Bigl(1-\frac1{2m}\Bigr)
  =1-\frac1{4m}-\frac{H_{2m}}{2},
\]
using $\sum_{k=1}^{2m-1}\frac1{k(k+1)}=1-\frac1{2m}$ and $H_{2m-1}=H_{2m}-\frac1{2m}$.

\subsection{The peak/valley bin \texorpdfstring{$\mathsf{onE}$}{onE}}
An on-foot position with an empty continuation side is a \emph{peak} (ascending block $\ell$
even, no unrevealed value above $v$) or a \emph{valley} (odd $\ell$, none below). At a peak,
$v$ is the maximum of block $\ell$ and every value above $v$ lies in a revealed block; that is,
$v=\max\{c:L_c\ge\ell\}$. Since the maximal-value card among those with label $\ge\ell$ has a
uniform label on $\{\ell,\dots,2m-1\}$, a peak occurs at block $\ell$ with probability
$1/(2m-\ell)$. There $G$ flips and guesses the maximum unrevealed value, which is correct iff
that value lies in block $\ell+1$ (whence it heads the next block); the maximal-value card among
labels $\ge\ell+1$ has uniform label, so this happens with probability $1/(2m-\ell-1)$. The
contribution is $\frac{1}{2m-\ell}\bigl(\frac{1}{2m-\ell-1}-\frac{1}{2m-\ell}\bigr)$. Valleys
mirror peaks. Writing $r=2m-\ell$, peaks (even $\ell$) give even $r\in\{2,\dots,2m\}$ and
valleys (odd $\ell$) give odd $r\in\{3,\dots,2m-1\}$ --- the case $r=1$ is the final output
card, which carries no guess --- so together $r$ ranges over $\{2,\dots,2m\}$ and
\[
  \mathsf{onE}=\sum_{r=2}^{2m}\Bigl(\frac{1}{r(r-1)}-\frac1{r^2}\Bigr)
  =\Bigl(1-\frac1{2m}\Bigr)-\bigl(H_{2m}^{(2)}-1\bigr)
  =2-\frac1{2m}-H_{2m}^{(2)}.
\]

\subsection{The vanishing bin \texorpdfstring{$\mathsf{offE}$}{offE}}
An off-foot position with an empty continuation side occurs mainly at the top block
($\ell=2m-1$ entered upward, nothing above): $G$ flips to the maximum unrevealed value below
$v$, which is the second card of $B_{2m-1}$ --- exactly the next card. There the hit is $1$ and
the reference is $1/(2m-\ell)=1$, so the excess is $0$; the remaining configurations have
vanishing probability. Hence $\mathsf{offE}\to0$.

\subsection{Assembly}
Summing \eqref{eq:pieces},
\[
  S(m)=\frac{H_{2m}-1}{2}+\Bigl(1-\frac1{4m}-\frac{H_{2m}}{2}\Bigr)+\Bigl(2-\frac1{2m}-H_{2m}^{(2)}\Bigr)
  =\frac52-\frac{3}{4m}-H_{2m}^{(2)},
\]
the $H_{2m}$ terms cancelling. With \eqref{eq:bS}, $b(m)=-1+\frac1{2m}+S(m)=\frac32-\frac1{4m}
-H_{2m}^{(2)}$. Together with Propositions~\ref{prop:slope} and \ref{prop:fade} this proves
Theorem~\ref{thm:main}. At $m=1$: $\mathsf{onC}=\tfrac14$, $\mathsf{offC}=0$,
$\mathsf{onE}=\tfrac14$, $\mathsf{offE}=0$, so $S(1)=\tfrac12$ and $b(1)=0$, i.e.\
$\Fg(n,1)=\tfrac34n$. \qed

\section{On the optimality of \texorpdfstring{$G$}{G}}\label{sec:opt}

Theorem~\ref{thm:main} is a statement about the explicit strategy $G$. Its relation to the
optimum is Clay's separate optimality claim. Three pieces of evidence bear on it.

\begin{enumerate}[topsep=2pt,itemsep=2pt]
  \item \emph{Parse-conditional optimality.} By Lemma~\ref{lem:key}, conditional on the true
        parse, $G$'s guess is the posterior-maximiser; so $G$ is optimal against an observer
        who additionally knows the block-parse. The only gap between $G$ and the true optimum
        is the parse uncertainty analysed in \S\ref{sec:intercept}.
  \item \emph{Exact finite verification.} A dynamic program over the sufficient statistic
        (direction, rank of the last card, ascending-run-length composition) computes $\eopt$
        in exact rational arithmetic. For every $n\le9$ and $m\le10$ one finds
        $\eopt(n,m)=\Fg(n,m)$ \emph{as exact fractions} (gap $0$) --- including cells with
        $n/2m\approx 0.4$, outside Clay's hedge.
  \item \emph{Deck-scale verification.} A common-random-numbers comparison of the greedy
        optimum against $G$, at $n\in\{52,104\}$ and $m=5,\dots,40$, finds
        $|\eopt-\Fg|\le0.01$ throughout (the only excursions are within sampling error and of
        the wrong sign for a genuine gap).
\end{enumerate}

We do \emph{not} prove $\eopt=\Fg$ in closed form for all $(n,m)$; like Clay we leave the
general optimality of $G$ as a conjecture, now with exact grid evidence and deck-scale
confirmation. Under it, \eqref{eq:main} is the optimal expected score, and Clay's
Conjecture~3 holds with the explicit constant and fade rate above. We emphasise that
Theorem~\ref{thm:main} itself --- the value of $G$ --- does not depend on this.

\section{Computational verification}\label{sec:verify}

Every identity above is machine-checked in exact arithmetic. (i) Lemma~\ref{lem:key} and the
hit formula \eqref{eq:hit} are verified against brute-force enumeration of all $(2m)^n$ label
vectors, matching as exact fractions for every prefix, parse and direction ($m=2,\ n\le7$;
$m=3,\ n\le6$). (ii) The two exchangeability identities behind $\mathsf{offC}$ and
$\mathsf{onE}$ --- that the maximal-value card among labels $\ge\ell$ has label $\ell$ with
probability $1/(2m-\ell)$, and the off-foot rate $\tfrac12-1/(2m+1-\ell)$ --- are verified
exactly by enumeration. (iii) The four-piece identity \eqref{eq:pieces} and the resulting
$b(m)$ are checked to reproduce the exact intercepts
\[
  b(1..8)=0,\ -\tfrac{7}{144},\ -\tfrac{269}{3600},\ -\tfrac{63449}{705600},\ -\tfrac{126713}{1270080},\ -\tfrac{16388909}{153679680},\ \dots
\]
obtained independently from an exact-rational dynamic program, whose intercept sequence is also
$C$-finite; its minimal recurrence factors over $\mathbb{Q}$ with characteristic roots
$\{i/m\}$, \emph{each of multiplicity $3$}, and $\{(2i-1)/2m\}$ (spectral gap $1/m$). The
dominant subdominant root $1-1/m$ and its multiplicity $3$ confirm, by a second route, both the
fade rate and the $n^2$ prefactor of Proposition~\ref{prop:fade}. (iv) The block-model bin
decomposition
converges to \eqref{eq:pieces} as $n$ grows (the per-bin convergence is slow,
$\bigoh((1-\tfrac1{2m})^n)$, as expected). The verification code is available from the author.

\section*{Acknowledgments}

This work grew out of a deterministic blackjack/shuffle simulator, from which the exact
posterior for shelf shuffles was reused to attack the guessing problem. The investigation ---
model, lemmas, the block decomposition, and the proofs above --- was carried out
collaboratively with a large language model (Anthropic's Claude); all mathematical claims are
independently verified by the exact-arithmetic enumeration described in \S\ref{sec:verify}. I
thank Alexander Clay for the paper that framed the problem.

\begin{thebibliography}{9}
\bibitem{DFH} P.~Diaconis, J.~Fulman and S.~Holmes,
  \emph{Analysis of casino shelf shuffling machines},
  Ann.\ Appl.\ Probab.\ \textbf{23} (2013), 1692--1720. arXiv:1107.2961.
\bibitem{Clay} A.~Clay,
  \emph{Guessing strategies for shelf-shuffling machines},
  arXiv:2507.10294 (2025).
\bibitem{Tripathi} (single-shelf follow-up, 2026), arXiv:2602.07920.
\bibitem{Kuba} (single-shelf follow-up, 2026), arXiv:2602.12928.
\bibitem{CKT} Clay, Kuba, Tripathi (2026), arXiv:2607.10418.
\end{thebibliography}

\end{document}
